\begin{table}[htbp]
\centering
\renewcommand{\arraystretch}{1.3}
\small
\begin{tabular}{ll ccc}
\toprule
\textbf{Métrica} & \textbf{Factor} & \textbf{Kruskal-Wallis H} & \textbf{valor p} & \textbf{$\eta^2$} \\
\midrule
Estabilidad & arquitectura & 91,1 & $<$0,0001 & 0,24 \\
Estabilidad & explicador & 125,6 & $<$0,0001 & 0,36 \\
Estabilidad & balanceo & 9,1 & 0,0107 & 0,01 \\
Estabilidad & escenario & 9,9 & 0,0423 & 0,04 \\
\midrule
Plaus. aristas & arquitectura & 10,5 & 0,0149 & 0,04 \\
Plaus. aristas & explicador & 195,8 & $<$0,0001 & 0,64 \\
Plaus. aristas & balanceo & 2,4 & 0,3081 & 0,01 \\
Plaus. aristas & escenario & 2,6 & 0,6268 & 0,01 \\
\midrule
Plaus. features & arquitectura & 11,3 & 0,0103 & 0,05 \\
Plaus. features & explicador & 176,2 & $<$0,0001 & 0,45 \\
Plaus. features & balanceo & 3,5 & 0,1730 & 0,01 \\
Plaus. features & escenario & 3,0 & 0,5660 & 0,01 \\
\bottomrule
\end{tabular}
\caption{Prueba de Kruskal-Wallis y tamaño de efecto de cada factor sobre cada métrica en el grafo sintético. El estadístico H corresponde a la prueba no paramétrica, que es la que decide la significancia, mientras que el $\eta^2$ es el tamaño de efecto del análisis de varianza sobre los mismos grupos, incluido como medida de magnitud comparable entre factores. El explicador ejerce el efecto dominante en las tres métricas, la arquitectura solo en la estabilidad, y tanto el balanceo como el escenario de desbalance resultan despreciables.}
\label{tab:synth-stats}
\end{table}
